
\section{另一个技术引理}
\label{sec:another-lemma}

我们现在来证明前面使用过的一条技术性引理\footnote{这一条比它的几个前身都要重要。}（我们曾用它来证明 $\widehat{\psi}$ 是一个\trImmersion）。

\emph{引理（\cref{thm:alpha-bounded-lp-limit}）。}序列 $(\alpha_{n})$ 是\trBounded{}的，并且
\[
\lim_{n\to+\infty}\Bigl\|\Bigl(\frac{\partial w_{\rho}}{\partial\rho}-\alpha_{n}\frac{\partial w_{\rho}}{\partial s}\Bigr)_{\rho_{n}}\Bigr\|_{L^{p}}=0\text{。}
\]

\emph{\cref{thm:alpha-bounded-lp-limit}的证明。}我们先来证明序列 $(\alpha_{n})$ 是\trBounded{}的。令
\[
A_{n}=\Bigl\|\Bigl(D_{1}\widetilde{\exp}_{(w_{\rho},\gamma(\rho))}\Bigl(\frac{\partial w_{\rho}}{\partial\rho}\Bigr)+D_{2}\widetilde{\exp}_{(w_{\rho},\gamma(\rho))}\Bigl(\frac{\partial\gamma}{\partial\rho}\Bigr)\Bigr)_{\rho=\rho_{n}}\Bigr\|_{L^{p}}\text{，}
\]
\[
B_{n}=\Bigl\|\Bigl(D_{1}\widetilde{\exp}_{(w_{\rho},\gamma(\rho))}\Bigl(\frac{\partial w_{\rho}}{\partial s}\Bigr)+D_{2}\widetilde{\exp}_{(w_{\rho},\gamma(\rho))}\Bigl(\frac{\partial\gamma}{\partial s}\Bigr)\Bigr)_{\rho=\rho_{n}}\Bigr\|_{L^{p}}\text{，}
\]
从而，按定义，我们有 $|\alpha_{n}|=A_{n}/B_{n}$。让我们来证明 $A_{n}$ 从上\trBounded，而 $B_{n}$ 从下被一个正常数控制。我们取定一个 $r_{0}>0$。我们有
\[
\lim_{\rho\to+\infty}\|\gamma\|_{W^{1,p}}=0\text{，}\qquad\text{从而 }\|\gamma(\rho_{n})\|_{W^{1,p}}\le r_{0}
\]
对充分大的 $n$ 成立，于是，再一次利用连续\trInjective
\[
W^{1,p}(\RR\times S^{1};\RR^{m})\hookrightarrow L^{\infty}(\RR\times S^{1};\RR^{m})\text{，}
\]
我们便有
\[
\sup_{(s,t)\in\RR\times S^{1}}\|\gamma(s,t)\|\le Kr_{0}\text{。}
\]
对一个取定的 $(s,t)$ 与 $\rho=\rho_{n}$，我们于是可以把\cref{thm:a-b-estimates} 应用于诸线性映射 $D_{i}\widetilde{\exp}_{(p,X)}(\cdot)$（$i=1,2$，取 $p=w_{\rho_{n}}(s,t)$ 与 $X=\gamma(\rho_{n})(s,t)$）。由此推出：对\trEuclidean{}范数，在 $(s,t)$ 处且 $\rho=\rho_{n}$ 时，我们有
\[
\Bigl\|D_{1}\widetilde{\exp}_{(w_{\rho},\gamma(\rho))}\Bigl(\frac{\partial w_{\rho}}{\partial\rho}\Bigr)\Bigr\|\le C\Bigl\|\frac{\partial w_{\rho}}{\partial\rho}\Bigr\|
\]
与\footnote{译注：其下不等式左端原书印作 $\Bigl\|D_{2}\widetilde{\exp}_{(w_{\rho},\gamma(\rho))}\Bigl(\dfrac{\partial w_{\rho}}{\partial\rho}\Bigr)\Bigr\|$；按 $A_{n}$ 的定义中 $D_{2}\widetilde{\exp}$ 的变元为 $\dfrac{\partial\gamma}{\partial\rho}$，且右端为 $C\Bigl\|\dfrac{\partial\gamma}{\partial\rho}\Bigr\|$，疑应作 $\Bigl(\dfrac{\partial\gamma}{\partial\rho}\Bigr)$。此处径改。——译者注}
\[
\Bigl\|D_{2}\widetilde{\exp}_{(w_{\rho},\gamma(\rho))}\Bigl(\frac{\partial\gamma}{\partial\rho}\Bigr)\Bigr\|\le C\Bigl\|\frac{\partial\gamma}{\partial\rho}\Bigr\|\text{。}
\]
取 $L^{p}$ 范数，我们于是有
\[
A_{n}\le C\Bigl(\Bigl\|\Bigl(\frac{\partial w_{\rho}}{\partial\rho}\Bigr)_{\rho=\rho_{n}}\Bigr\|_{L^{p}}+\Bigl\|\Bigl(\frac{\partial\gamma}{\partial\rho}\Bigr)_{\rho=\rho_{n}}\Bigr\|_{L^{p}}\Bigr)\text{。}
\]
回忆一下
\[
w_{\rho}(s,t)\left\{\begin{array}{ll}
=u(s+\rho,t) & \text{对 } s\le-1\\
\text{趋于 } y(t) & \text{对 } s\in[-1,1]\text{ 且 }\rho\to+\infty\\
=v(s-\rho,t) & \text{对 } s\ge 1\text{，}
\end{array}\right.
\]
从而
\[
\Bigl\|\frac{\partial w_{\rho}}{\partial\rho}\Bigr\|_{L^{p}}\le\Bigl\|\frac{\partial u}{\partial s}\Bigr\|_{L^{p}}+\Bigl\|\frac{\partial v}{\partial s}\Bigr\|_{L^{p}}+1
\]
对充分大的 $\rho$ 成立。特别地，这个范数关于 $\rho$ 一致\trBounded。我们还有
\[
\lim_{\rho\to+\infty}\Bigl\|\frac{\partial\gamma}{\partial\rho}\Bigr\|=0
\]
（由\cref{thm:dgamma-rho-tends-zero}）\footnote{译注：上式极限号下的变量原书印作 $s\to+\infty$；按其左端只依赖于 $\rho$，且所援引的\cref{thm:dgamma-rho-tends-zero}断言的正是 $\lim_{\rho\to+\infty}\Bigl\|\dfrac{\partial\gamma}{\partial\rho}\Bigr\|_{W^{1,p}}=0$，疑应为 $\rho\to+\infty$。此处径改。——译者注}。于是，序列 $A_{n}$ 是\trBounded{}的。现在我们转向 $B_{n}$。对 $\rho=\rho_{n}$，我们有
\[
\begin{aligned}
B_{n}&\ge\Bigl\|D_{1}\widetilde{\exp}_{(w_{\rho},\gamma(\rho))}\frac{\partial w_{\rho}}{\partial s}\Bigr\|_{L^{p}}-\Bigl\|D_{2}\widetilde{\exp}_{(w_{\rho},\gamma(\rho))}\frac{\partial\gamma}{\partial s}\Bigr\|_{L^{p}}\\
&\ge\Bigl\|D_{1}\widetilde{\exp}_{(w_{\rho},0)}\frac{\partial w_{\rho}}{\partial s}\Bigr\|_{L^{p}}-\Bigl\|D_{1}\widetilde{\exp}_{(w_{\rho},\gamma(\rho))}\frac{\partial w_{\rho}}{\partial s}-D_{1}\widetilde{\exp}_{(w_{\rho},0)}\frac{\partial w_{\rho}}{\partial s}\Bigr\|_{L^{p}}\\
&\qquad\qquad-C\Bigl\|\frac{\partial\gamma}{\partial s}\Bigr\|_{L^{p}}
\end{aligned}
\]
（对第三项，用的是上面 $D_{2}\widetilde{\exp}$ 下界估计在 $\partial s$ 情形的类比式）。我们取定一个对子 $(s,t)\in\RR\times S^{1}$ 并（再一次）令 $\rho=\rho_{n}$。由\cref{thm:a-b-estimates}，我们有
\[
\begin{aligned}
\Bigl\|D_{1}\widetilde{\exp}_{(w_{\rho},\gamma(\rho))}\frac{\partial w_{\rho}}{\partial s}-D_{1}\widetilde{\exp}_{(w_{\rho},0)}\frac{\partial w_{\rho}}{\partial s}\Bigr\|&\le C\,\|\gamma(\rho)\|\,\Bigl\|\frac{\partial w_{\rho}}{\partial s}\Bigr\|\\
&\le Ck\,\|\gamma(\rho)\|\quad\text{因为 }\Bigl\|\frac{\partial w_{\rho}}{\partial s}\Bigr\|_{\cC^{0}}\le k\text{。}
\end{aligned}
\]
利用等式 $D_{1}\widetilde{\exp}_{(w_{\rho},0)}=\operatorname{Id}$，我们便得到
\[
B_{n}\ge\Bigl\|\frac{\partial w_{\rho}}{\partial s}\Bigr\|_{L^{p}}-Ck\,\|\gamma(\rho)\|_{L^{p}}-C\Bigl\|\frac{\partial\gamma}{\partial s}\Bigr\|_{L^{p}}\text{。}
\]
由于
\[
\frac{\partial w_{\rho}}{\partial s}(s,t)=\frac{\partial u}{\partial s}(s+\rho,t)\qquad\text{对 } s\le-1\text{，}
\]
我们有\footnote{译注：上式条件原书印作“对 $s\ge-1$”；按本页开头 $w_{\rho}$ 的分段定义（$w_{\rho}(s,t)=u(s+\rho,t)$ 当 $s\le-1$）以及紧随其后的积分区域 $]-\infty,\rho-1]\times S^{1}$，疑应为 $s\le-1$。此处径改。——译者注}
\[
\Bigl\|\frac{\partial w_{\rho}}{\partial s}\Bigr\|_{L^{p}}\ge\Bigl\|\frac{\partial u}{\partial s}\Bigr\|_{L^{p}(]-\infty,\rho-1]\times S^{1})}\ge C_{0}>0
\]
而（由\cref{thm:gamma-rho-tends-zero}）
\[
\lim_{\rho\to+\infty}\|\gamma(\rho)\|_{L^{p}}=0\text{，}\qquad\lim_{\rho\to+\infty}\Bigl\|\frac{\partial\gamma}{\partial s}\Bigr\|_{L^{p}}=\lim_{\rho\to+\infty}\|\gamma\|_{W^{1,p}}=0\text{。}
\]
于是存在一个常数 $C_{1}$，使得对每个 $n$ 都有 $B_{n}\ge C_{1}$。从而，序列 $(\alpha_{n})$ 是\trBounded{}的。

其次，按假设，我们有
\[
\frac{\partial}{\partial\rho}\Bigl(\exp_{w_{\rho}}\gamma_{\rho}(s,t)\Bigr)_{\rho_{n}}=\alpha_{n}\frac{\partial}{\partial s}\Bigl(\exp_{w_{\rho}}\gamma_{\rho}(s,t)\Bigr)_{\rho_{n}}\text{，}
\]
亦即
\[
\frac{\partial}{\partial\rho}\Bigl(\widetilde{\exp}_{w_{\rho}}\gamma_{\rho}(s,t)\Bigr)_{\rho_{n}}=\alpha_{n}\frac{\partial}{\partial s}\Bigl(\widetilde{\exp}_{w_{\rho}}\gamma_{\rho}(s,t)\Bigr)_{\rho_{n}}\text{，}
\]
也就是说，
\[
\begin{aligned}
&\Bigl(D_{1}\widetilde{\exp}_{(w_{\rho},\gamma(\rho))}\frac{\partial w_{\rho}}{\partial\rho}+D_{2}\widetilde{\exp}_{(w_{\rho},\gamma(\rho))}\frac{\partial\gamma}{\partial\rho}\Bigr)_{\rho_{n}}\\
&\quad=\alpha_{n}\Bigl(D_{1}\widetilde{\exp}_{(w_{\rho},\gamma(\rho))}\frac{\partial w_{\rho}}{\partial s}+D_{2}\widetilde{\exp}_{(w_{\rho},\gamma(\rho))}\frac{\partial\gamma}{\partial s}\Bigr)_{\rho_{n}}\text{。}
\end{aligned}
\]
上面用过的那些估计表明
\[
\lim_{n\to+\infty}\Bigl\|D_{2}\Bigl(\widetilde{\exp}_{(w_{\rho},\gamma(\rho))}\frac{\partial\gamma}{\partial\rho}\Bigr)_{\rho_{n}}\Bigr\|_{L^{p}}=\lim_{n\to+\infty}\Bigl\|D_{2}\Bigl(\widetilde{\exp}_{(w_{\rho},\gamma(\rho))}\frac{\partial\gamma}{\partial s}\Bigr)_{\rho_{n}}\Bigr\|_{L^{p}}=0\text{。}
\]
由于序列 $\alpha_{n}$ 是\trBounded{}的，我们有
\[
\lim_{n\to+\infty}\Bigl\|\Bigl(D_{1}\widetilde{\exp}_{(w_{\rho},\gamma(\rho))}\Bigl(\frac{\partial w_{\rho}}{\partial\rho}-\alpha_{n}\frac{\partial w_{\rho}}{\partial s}\Bigr)\Bigr)_{\rho_{n}}\Bigr\|_{L^{p}}=0\text{。}
\]
我们再一次取定一个对子 $(s,t)\in\RR\times S^{1}$。由\cref{thm:a-b-estimates}，对\trEuclidean{}范数，在 $(s,t)$ 处且 $\rho=\rho_{n}$ 时，我们有
\[
\begin{aligned}
&\Bigl\|D_{1}\widetilde{\exp}_{(w_{\rho},\gamma(\rho))}\Bigl(\frac{\partial w_{\rho}}{\partial\rho}-\alpha_{n}\frac{\partial w_{\rho}}{\partial s}\Bigr)-D_{1}\widetilde{\exp}_{(w_{\rho},0)}\Bigl(\frac{\partial w_{\rho}}{\partial\rho}-\alpha_{n}\frac{\partial w_{\rho}}{\partial s}\Bigr)\Bigr\|\\
&\qquad\qquad\le C\,\|\gamma(\rho)\|\,\Bigl\|\frac{\partial w_{\rho}}{\partial\rho}-\alpha_{n}\frac{\partial w_{\rho}}{\partial s}\Bigr\|\\
&\qquad\qquad\qquad\le Ck\,\|\gamma(\rho)\|
\end{aligned}
\]
因为 $\Bigl\|\dfrac{\partial w_{\rho}}{\partial\rho}\Bigr\|_{\cC^{0}}$、$\Bigl\|\dfrac{\partial w_{\rho}}{\partial s}\Bigr\|_{\cC^{0}}$ 与 $\alpha_{n}$ 都是\trBounded{}的。

取 $L^{p}$ 范数并利用 $\lim_{\rho\to+\infty}\|\gamma\|_{L^{p}}=0$ 这一事实，我们推出：表达式
\[
D_{1}\widetilde{\exp}_{(w_{\rho},\gamma(\rho))}\Bigl(\frac{\partial w_{\rho}}{\partial\rho}-\alpha_{n}\frac{\partial w_{\rho}}{\partial s}\Bigr)-D_{1}\widetilde{\exp}_{(w_{\rho},0)}\Bigl(\frac{\partial w_{\rho}}{\partial\rho}-\alpha_{n}\frac{\partial w_{\rho}}{\partial s}\Bigr)
\]
的 $L^{p}$ 范数（取 $\rho=\rho_{n}$）当 $n$ 趋于无穷时趋于 $0$。于是，既然我们已经看到第一项的范数趋于 $0$，
\[
\lim_{n\to+\infty}\Bigl\|D_{1}\widetilde{\exp}_{(w_{\rho},0)}\Bigl(\frac{\partial w_{\rho}}{\partial\rho}-\alpha_{n}\frac{\partial w_{\rho}}{\partial s}\Bigr)\Bigr\|_{L^{p}}=0\text{，}
\]
也就是说，
\[
\lim_{n\to+\infty}\Bigl\|\Bigl(\frac{\partial w_{\rho}}{\partial\rho}-\alpha_{n}\frac{\partial w_{\rho}}{\partial s}\Bigr)_{\rho_{n}}\Bigr\|_{L^{p}}=0\text{，}
\]
这正是我们所要证明的。\noqed
